\begin{table}[H]
  \centering
  \caption{Caso de uso CUN-01: Gestionar ventas (registrar venta)}
  \label{tab:cun_01}
  {\small
  \setlength{\tabcolsep}{4pt}
  \renewcommand{\arraystretch}{1.05}
  \begin{tabularx}{\textwidth}{|>{\bfseries\raggedright\arraybackslash}p{3.2cm}|>{\raggedright\arraybackslash}X|}
    \hline
    Caso & Gestionar ventas (registrar venta) \\
    \hline
    Descripción & Permite al vendedor procesar la salida de productos solicitados por el cliente dentro del sistema de inventario. \\
    \hline
    Precondición & El usuario (vendedor o cajero) debe haber iniciado sesión y el cliente debe haber seleccionado los juguetes a comprar con stock disponible. \\
    \hline
    Secuencia &
    1. El vendedor ingresa la opción de registrar venta en el sistema. \newline
    2. El vendedor selecciona o ingresa los productos solicitados por el cliente. \newline
    3. El sistema valida el stock disponible de los juguetes. \newline
    4. El vendedor confirma la transacción de venta. \newline
    5. El sistema procesa el registro de la venta y ejecuta de manera automática la actualización del inventario. \newline
    6. El sistema muestra un mensaje de confirmación de venta exitosa. \\
    \hline
    Errores/Flujos alternativos & Stock insuficiente para la cantidad solicitada, error en la conexión o fallo al procesar el registro. \\
    \hline
    Postcondición & La venta queda registrada en el sistema y el inventario de los juguetes vendidos se actualiza automáticamente. \\
    \hline
    Objetivo & Agilizar el proceso de venta y mantener el nivel de stock actualizado en tiempo real en la tienda. \\
    \hline
  \end{tabularx}}
  \par\vspace{0.2ex}
  \raggedright\noindent\small \emph{Nota.} Elaboración propia (2026). Cuadro de especificación fundamental del CUN-01.
\end{table}

\begin{table}[H]
  \centering
  \caption{Caso de uso CUN-02: Gestionar abastecimiento}
  \label{tab:cun_02}
  {\small
  \setlength{\tabcolsep}{4pt}
  \renewcommand{\arraystretch}{1.05}
  \begin{tabularx}{\textwidth}{|>{\bfseries\raggedright\arraybackslash}p{3.2cm}|>{\raggedright\arraybackslash}X|}
    \hline
    Caso & Gestionar abastecimiento \\
    \hline
    Descripción & El proceso de cotizar, hacer el pedido al proveedor y recepcionar los nuevos juguetes. \\
    \hline
    Precondición & El administrador identifica la necesidad de reponer stock en la tienda para evitar quiebres de inventario. \\
    \hline
    Secuencia &
    1. El administrador revisa el catálogo de productos y detecta baja disponibilidad o la necesidad de nuevos artículos. \newline
    2. El administrador cotiza los productos con el proveedor externo. \newline
    3. El administrador realiza y emite el pedido formal al proveedor. \newline
    4. El proveedor despacha y envía la nueva mercancía solicitada hacia la tienda. \newline
    5. El administrador o encargado de almacén recibe y recepciona los nuevos juguetes en el establecimiento. \newline
    6. El sistema o el administrador registra la entrada de la nueva mercancía para actualizar los niveles de stock. \\
    \hline
    Errores/Flujos alternativos & Falta de stock en el proveedor, retraso en la entrega de la mercadería o discrepancias en la recepción de los juguetes. \\
    \hline
    Postcondición & El inventario de la tienda es abastecido con los nuevos productos y los niveles de stock quedan actualizados. \\
    \hline
    Objetivo & Mantener el stock de la tienda actualizado y evitar quiebres de inventario mediante una correcta provisión de mercancía. \\
    \hline
  \end{tabularx}}
  \par\vspace{0.2ex}
  \raggedright\noindent\small \emph{Nota.} Elaboración propia (2026). Cuadro de especificación fundamental del CUN-02.
\end{table}

\begin{table}[H]
  \centering
  \caption{Caso de uso CUN-03: Controlar inventario}
  \label{tab:cun_03}
  {\small
  \setlength{\tabcolsep}{4pt}
  \renewcommand{\arraystretch}{1.05}
  \begin{tabularx}{\textwidth}{|>{\bfseries\raggedright\arraybackslash}p{3.2cm}|>{\raggedright\arraybackslash}X|}
    \hline
    Caso & Controlar inventario \\
    \hline
    Descripción & El proceso interno de realizar auditorías, registrar mermas (juguetes rotos) y actualizar los niveles mínimos de stock. \\
    \hline
    Precondición & El administrador o personal autorizado requiere verificar físicamente la mercancía y el estado del almacén. \\
    \hline
    Secuencia &
    1. El administrador accede al módulo de control de inventario. \newline
    2. Se realiza una auditoría física de los juguetes almacenados en la tienda. \newline
    3. Se identifican y registran las mermas o artículos dañados. \newline
    4. Se actualizan los niveles mínimos de stock y los registros de existencias. \newline
    5. El sistema guarda los cambios y emite un informe actualizado de inventario. \\
    \hline
    Errores/Flujos alternativos & Discrepancias entre el stock físico y el sistema, o registro de mermas con datos erróneos. \\
    \hline
    Postcondición & El inventario físico y del sistema coinciden, con las mermas registradas y saldos actualizados. \\
    \hline
    Objetivo & Mantener la precisión de existencias, prevenir pérdidas económicas y asegurar control interno. \\
    \hline
  \end{tabularx}}
  \par\vspace{0.2ex}
  \raggedright\noindent\small \emph{Nota.} Elaboración propia (2026). Cuadro de especificación fundamental del CUN-03.
\end{table}

\begin{table}[H]
  \centering
  \caption{Caso de uso CUN-04: Gestionar devoluciones y cambios}
  \label{tab:cun_04}
  {\small
  \setlength{\tabcolsep}{4pt}
  \renewcommand{\arraystretch}{1.05}
  \begin{tabularx}{\textwidth}{|>{\bfseries\raggedright\arraybackslash}p{3.2cm}|>{\raggedright\arraybackslash}X|}
    \hline
    Caso & Gestionar devoluciones y cambios \\
    \hline
    Descripción & El proceso desde que un cliente reporta un juguete defectuoso hasta que se resuelve con cambio o reembolso. \\
    \hline
    Precondición & El cliente se presenta en el establecimiento con el juguete y el comprobante de pago original. \\
    \hline
    Secuencia &
    1. El cliente se acerca al vendedor para reportar un juguete defectuoso o solicitar cambio. \newline
    2. El vendedor recibe el producto y solicita el comprobante de pago. \newline
    3. El vendedor verifica el estado del juguete y las condiciones de la devolución. \newline
    4. El vendedor procesa la solicitud seleccionando cambio de producto o reembolso. \newline
    5. El sistema actualiza el registro de inventario o genera la nota de crédito respectiva. \newline
    6. Se entrega el nuevo producto o se efectúa el reembolso al cliente. \\
    \hline
    Errores/Flujos alternativos & Comprobante inválido, producto fuera del plazo permitido o daños atribuibles a mal uso. \\
    \hline
    Postcondición & La devolución queda registrada y el inventario se ajusta según la solución aplicada. \\
    \hline
    Objetivo & Atender eficientemente incidencias de clientes y garantizar la satisfacción y trazabilidad. \\
    \hline
  \end{tabularx}}
  \par\vspace{0.2ex}
  \raggedright\noindent\small \emph{Nota.} Elaboración propia (2026). Cuadro de especificación fundamental del CUN-04.
\end{table}

\begin{table}[H]
  \centering
  \caption{Caso de uso CUN-05: Flujo de venta y reposici\'on (trabajadores del negocio)}
  \label{tab:cun_05}
  {\small
  \setlength{\tabcolsep}{4pt}
  \renewcommand{\arraystretch}{1.05}
  \begin{tabularx}{\textwidth}{|>{\bfseries\raggedright\arraybackslash}p{3.2cm}|>{\raggedright\arraybackslash}X|}
    \hline
    Caso & Flujo de venta y reposición (trabajadores del negocio) \\
    \hline
    Descripción & Describe la interacción interna de los empleados y el sistema desde la venta hasta la reposición de mercancía. \\
    \hline
    Precondición & Existen productos registrados y la tienda se encuentra formalmente abierta al público. \\
    \hline
    Secuencia &
    1. El Vendedor registra la venta de los productos solicitados. \newline
    2. El Sistema actualiza el stock provisional. \newline
    3. El Cajero cobra al cliente y emite el comprobante de pago. \newline
    4. El Administrador verifica niveles de stock y emite orden de reposición. \newline
    5. El Encargado de almacén recibe la mercancía física y valida los ítems. \newline
    6. El Sistema actualiza definitivamente el inventario en bodega. \\
    \hline
    Errores/Flujos alternativos & Fallas de red interna, desabastecimiento en mostrador o retrasos en la estiba de almacén. \\
    \hline
    Postcondición & La venta queda saldada, el stock sincronizado y el almacén reabastecido. \\
    \hline
  \end{tabularx}}
  \par\vspace{0.2ex}
  \raggedright\noindent\small \emph{Nota.} Elaboración propia (2026). Cuadro de especificación del flujo interno (CUN-05).
\end{table}

\begin{table}[H]
  \centering
  \caption{Caso de uso CUN-06: Flujo con actores externos de negocio}
  \label{tab:cun_06}
  {\small
  \setlength{\tabcolsep}{4pt}
  \renewcommand{\arraystretch}{1.05}
  \begin{tabularx}{\textwidth}{|>{\bfseries\raggedright\arraybackslash}p{3.2cm}|>{\raggedright\arraybackslash}X|}
    \hline
    Caso & Flujo con actores externos de negocio \\
    \hline
    Descripción & Muestra la interacción entre Clientes y Proveedores externos con los puntos de contacto de la juguetería. \\
    \hline
    Precondición & El cliente solicita una compra o la empresa emite una orden de compra hacia el proveedor mayorista. \\
    \hline
    Secuencia &
    1. El Cliente solicita la compra de juguetes en la tienda. \newline
    2. El Vendedor procesa la solicitud y el Cajero efectúa el cobro. \newline
    3. El Cliente recibe su compra embalada con su comprobante legal. \newline
    4. El Administrador transmite la orden de reposición al Proveedor. \newline
    5. El Proveedor despacha el lote de juguetes hacia la tienda. \newline
    6. El Sistema registra formalmente el ingreso de la mercancía. \\
    \hline
    Errores/Flujos alternativos & Rechazo del pago por parte del cliente, falta de existencias en el proveedor o demoras logísticas. \\
    \hline
    Postcondición & El cliente recibe su producto y la tienda completa el aprovisionamiento. \\
    \hline
  \end{tabularx}}
  \par\vspace{0.2ex}
  \raggedright\noindent\small \emph{Nota.} Elaboración propia (2026). Cuadro de especificación del flujo externo (CUN-06).
\end{table}
